\subsubsection{Select Worked Problems}

\begin{problem}
How many ways can you arrange the letters in the word BANANA?
\end{problem}

\begin{solution}
There are $n = 6$ elements. Repeated items have frequency 3 and 2 (Ns and As). Then the number of arrangements are $\frac{6!}{3! \times 2!} = \frac{720}{12} = 60$.
\end{solution}

\medskip

\begin{problem}
You roll 2 dice. How many distinct outcomes are there? What if you can't tell the dice apart?
\end{problem}

\begin{solution}
With two dice, there are $6 \times 6 = 36$ outcomes if we can tell the dice apart. If we \textit{can't} tell the dice apart, then we can solve either logically or formulaically:

Logically, we can see that from the original 36 combinations, there are 6 outcomes where the dice are the same, leaving 30 outcomes using mixed pairs of dice. Of these, half are duplicate when we can't tell the dice apart, leaving 15 outcomes. We can then add the 6 double outcomes back in for a total of 21 outcomes.

Formulaically, we can solve using a similar approach, taking the sum of mixed pair outcomes and double outcomes as $C(6, 2) + C(6, 1) = 15 + 6 = 21$.

Lastly, we can consider this problem using the stars and bars approach, where we notice that there are $n=2$ identical dice and $6$ distinct bins with $k=5$ dividers between them. Then we can count the outcomes by $C(2 + 6 - 1, 6 -1) = C(7, 5) = 21$
\end{solution}

\medskip

\begin{problem}
How many ways can you distribute 8 identical balls into 3 distinct boxes?
\end{problem}

\begin{solution}
This is a direct application of stars and bars with $n=8$ and $k=3$. The number of outcomes is $C(8+3-1, 3-1) = C(10,2) = \frac{10 \times 9 \times 8 \times \cdots \times 3}{2} = 45$.
\end{solution}

\medskip

\begin{problem}
A deck has 52 cards. How many 5-card hands contain at least one ace?
\end{problem}

\begin{solution}
Since this solution asks for the number of hands containing at least one ace, it will be simpler to compute the complement. This is the number of hands containing no aces, we we then subtract from the total number of 5 card hands: $C(52, 5) - C(48, 5) = 886,656$.
\end{solution}

\medskip

\begin{problem}
You have 5 books on a shelf. In how many ways can you choose 3 books and arrange them left-to-right on a new shelf?
\end{problem}

\begin{solution}
This is simply $C(5,3) \times P(3, 3) = \frac{5 \times 4 \times 3}{3!} \times \frac{3!}{0!} = \frac{60 \times 3!}{3!} = 60$.
\end{solution}

\medskip

\begin{problem}
A pizza place offers 8 toppings. How many different pizzas can you make if you can choose any number of toppings (including none)?
\end{problem}

\begin{solution}
We can solve this by taking the summation of $C(8, i)$ from $i=0$ to $i=8$:
\[
Total\ outcomes = \sum\limits_{i=0}^{8} C(8,i) = 256
\]

There is a more succinct solution we can find if we approach the problem seeing each topping as the binary choice of being either on or off of the pizza, then the total outcomes are $2^8 = 256$.
\end{solution}

\medskip

\begin{problem}
How many ways can 12 people be split into 3 groups of 4?
\end{problem}

\begin{solution}
We choose each group independently with the remaining members of the pool, and then divide by the reorderings since the groups are indistinguishable:
\[
\frac{C(12, 4) \times C(8, 4) \times C(4, 4)}{3!} = \frac{495 \times 70 \times 1}{6} = 34,650.
\]
\end{solution}
